\chapter*{Resumo}
\paragraph{}Este relatório apresenta a conceção e implementação do Hubly, uma aplicação web desenvolvida no âmbito da unidade curricular de Projeto e Seminário. O projeto aborda a falta de centralização da informação nos processos de marketing de influência, nos quais empresas e criadores de conteúdo digital recorrem frequentemente a redes sociais, plataformas de mensagens e canais de email distintos para descobrir contactos, negociar colaborações e acompanhar o estado de potenciais parcerias. Esta fragmentação dificulta a comparação de oportunidades, a preservação do contexto de comunicação e a gestão consistente de múltiplos processos de contacto.

A solução proposta, Hubly, pretende responder a esta limitação através da disponibilização de um ambiente centralizado onde os utilizadores se podem registar como empresas ou como criadores. As contas de criadores são organizadas através de perfis sociais específicos por plataforma, como, por exemplo, uma conta de YouTube ou Instagram, permitindo que o mesmo criador possa gerir vários perfis de acordo com a sua presença real nas diferentes plataformas. Cada canal social é representado de forma independente, com a sua própria informação de audiência, setores de conteúdo e condições comerciais.

Empresas e criadores podem procurar-se mutuamente utilizando critérios como setor, plataforma, dimensão da audiência, intervalo de preços, dimensão da empresa e país, tornando o processo de descoberta mais estruturado e alinhado com os objetivos de cada utilizador. Além disso, o Hubly disponibiliza conversas internas associadas à empresa ou ao perfil social do criador relevante, garantindo que a comunicação permanece ligada ao contexto de uma potencial parceria.

Para responder a necessidades organizacionais, o Hubly introduz a gestão colaborativa de equipas, permitindo que vários utilizadores da mesma organização possam gerir e interagir com a plataforma através de uma conta corporativa partilhada. Adicionalmente, para assegurar transparência e responsabilidade no sistema, foi implementado um mecanismo de registo de auditoria, que permite manter um histórico de operações críticas realizadas na aplicação.

O protótipo resultante disponibiliza um ambiente estruturado para a gestão de interações entre empresas e criadores, reduzindo a dependência de canais externos e facilitando processos de colaboração mais organizados, escaláveis e transparentes.